\begin{table*}[t]
\centering
\scriptsize
\renewcommand{\arraystretch}{0.88}
\begin{tabularx}{\linewidth}{>{\raggedright\arraybackslash}p{0.17\linewidth}>{\raggedright\arraybackslash}p{0.21\linewidth}>{\raggedright\arraybackslash}p{0.22\linewidth}>{\raggedright\arraybackslash}p{0.16\linewidth}X}
\toprule
\textbf{Comparison} & \textbf{Why it looks comparable} & \textbf{Hidden mismatch} & \textbf{Metric pitfall} & \textbf{Normalised reading} \\
\midrule
PoisonedRAG vs. dense-retriever poisoning & Both insert malicious evidence into a retrieval corpus. & One evaluates end-to-end answer steering; the other may evaluate retriever ranking before generation. & Retrieval success and answer attack success are not the same target. & Compare only after recording whether poisoned evidence is retrieved and whether the generator changes the final answer. \\
Certified RAG defense vs. prompt-injection defense & Both are RAG defenses against corrupted retrieved content. & Certification assumes a bounded number of corrupted passages; prompt injection exploits instruction following rather than only false facts. & A certificate under bounded corruption may not imply instruction-hijack resistance. & Treat certification, instruction isolation, and document sanitisation as complementary defenses. \\
Traceback vs. prevention & Both reduce harm after poisoning. & Traceback identifies influential poisoned documents after failure; prevention attempts to block harmful output before response. & Lower post-hoc traceback error does not equal lower user-facing attack success. & Report prevention, detection, demotion, traceback, and recovery as separate outcome variables. \\
Backdoor trigger vs. natural query attack & Both can produce targeted malicious behavior. & Trigger attacks may rely on rare phrases; natural-query attacks must survive ordinary retrieval and ranking. & Attack success under attacker-crafted triggers overstates risk under natural user queries. & Require query-control and trigger-frequency fields before comparing attack strength. \\
\bottomrule
\end{tabularx}
\caption{Concrete examples of threat-model normalisation. These cases show how reported attack or defense success can reverse when attacker access, retrieval realism, and evaluation target differ.}
\label{tab:threat-model-examples}
\end{table*}

